% !PRINTED 577
\weberpage{577}
%<<CONT
und daraus schliesst man wieder nach Satz 3. (auf $\mathfrak p_1$ angewandt)
\[ \Phi_1(t) \equiv P_1(t) \Phi_2(t) \pmod p. \]
Hierin lässt sich dieselbe Betrachtung wiederholen, die zu der Congruenz
\[ \Phi_2(t) \equiv 0 \pmod{\mathfrak p_2 \ldots} \]
führt, woraus wieder nach 3.
\[ \Phi_2(t) \equiv P_2(t) \Phi_3(t) \pmod p \]
zu schliessen ist. Fährt man damit fort, bis alle Primfactoren von $p$ berücksichtigt sind, so ergiebt sich der folgende Satz:

5. \gesperrt{Ist $\Phi(t)$ eine ganze Function in $R$, die durch die Substitution $t=\tau$ durch $p$ theilbar wird, so lässt sich eine andere ganze Function $\Phi_0(t)$ in $R$ so bestimmen, dass}
\weq{29}{\Phi(t) \equiv \Phi_0(t) P(t) P_1(t) P_2(t) \ldots \pmod p}
\gesperrt{wird.}

\gesperrt{Das Product $P(t) P_1(t) P_2(t) \ldots$ ist vom Grade $n = f+f_1+f_2+\cdots$ und ist die ganze rationale Form niedrigsten Grades, die durch die Substitution $t=\tau$ durch $p$ theilbar wird.}

Zu den im Satze 5. vorkommenden Functionen $\Phi(t)$ gehört auch die ganze Function $n^{\mathrm{ten}}$ Grades
\[ F(t) = N(t-\tau), \]
die für $t=\tau$ verschwindet. Für diese Function wird $\Phi_0(t)$ von $t$ unabhängig, und da sowohl in $F(t)$ als in $P(t), P_1(t), P_2(t) \ldots$ die höchste Potenz von $t$ den Coëfficienten 1 hat, so wird $\Phi_0(t) = 1$. Es gilt also die Congruenz
\weq{30}{N(t-\tau) \equiv P(t) P_1(t) P_2(t) \ldots \pmod p.}
Fassen wir in der Zerlegung (25) der Primzahl $p$ die gleichen Primfactoren zu Potenzen zusammen, so können wir, wenn $\mathfrak p_1, \mathfrak p_2, \ldots$ verschiedene Primideale der Grade $f_1, f_2, \ldots$ sind, die positiven Exponenten $e_1, e_2, \ldots$ so annehmen, dass
\weq{31}{p = \mathfrak p_1^{e_1} \mathfrak p_2^{e_2} \ldots,}
und
\weq{32}{n = e_1 f_1 + e_2 f_2 + \cdots}
wird. Die Formel (30) nimmt dann die Gestalt an
\weq{33}{N(t-\tau) \equiv [P_1(t)]^{e_1} [P_2(t)]^{e_2} \ldots \pmod p.}
